\section{Fundamentos teóricos}\label{sec:theory}

\subsection{Balance de energía superficial y METRIC}

En el instante del paso del satélite, el balance de energía de la superficie
se escribe
\begin{equation}
  \Rn = G + H + LE \quad\Longrightarrow\quad LE = \Rn - G - H,
  \label{eq:balance}
\end{equation}
con todos los flujos en \si{\watt\per\square\metre}. La radiación neta es
\begin{equation}
  \Rn = (1-\alpha)\,R_{s\downarrow} + R_{L\downarrow} - R_{L\uparrow}
        - (1-\varepsilon_0)\,R_{L\downarrow},
  \label{eq:rn}
\end{equation}
donde $\alpha$ es el albedo de banda ancha, $R_{L\uparrow} =
\varepsilon_0\sigma T_s^4$ la radiación emitida por la superficie (MODIS
provee $T_s$ y permite estimar $\varepsilon_0$) y $R_{L\downarrow} =
\varepsilon_a\sigma T_a^4$ la radiación de onda larga incidente, con
$\varepsilon_a = 0{,}85(-\ln\tau_{sw})^{0{,}09}$ \parencite{waters2002sebal}.
La transmisividad de onda corta $\tau_{sw}$ depende de la presión
atmosférica y del agua precipitable (calculada a partir de la presión de
vapor $e_a$) en la forma particionada de Allen \parencite{asce2005standardized}.
En EarthFlux, $T_a$ y $e_a$ son, por tanto, las primeras dos vías de entrada
del forzamiento meteorológico.

El flujo de calor sensible se expresa como
\begin{equation}
  H = \rho_a c_p \frac{dT}{r_{ah}}, \qquad dT = a + b\,T_{s,datum},
  \label{eq:h}
\end{equation}
con $r_{ah}$ la resistencia aerodinámica al transporte de calor entre
\SI{0,1}{\metre} y \SI{2}{\metre} sobre el plano de desplazamiento, que
depende de la velocidad de fricción $u_*$ y, a través de ella, del viento
extrapolado logarítmicamente a una altura de mezcla de \SI{200}{\metre}
($u_{200}$; tercera vía de entrada). Los coeficientes $a$ y $b$ se obtienen
imponiendo en el ancla caliente $H = \Rn - G - LE_{hot}$ con $LE_{hot}=0$ y
en el ancla fría $LE_{cold} = 1{,}05\,\lambda\,\ETr$ (cuarta vía de entrada:
$\ETr$ instantánea en el ancla fría), iterando la corrección de estabilidad
de Monin--Obukhov hasta la convergencia \parencite{allen2007metric1,waters2002sebal}.
Con $LE$ como residuo, la evapotranspiración instantánea es $ET_{inst} =
3600\,LE/(\lambda\rho_w)$ y la diaria
\begin{equation}
  \ETa = \ETrF \cdot \ETr^{24}, \qquad \ETrF = \frac{ET_{inst}}{\ETr},
  \label{eq:eta}
\end{equation}
que introduce la quinta vía de entrada, la $\ETr$ de 24~h. La calibración
interna transforma un sesgo uniforme de $T_a$ o de $u$ en un cambio de $a$ y
$b$; un sesgo que \emph{varía} en el espacio, en cambio, no puede absorberse
con dos parámetros. Esa asimetría motiva el diseño de sustitución de una
variable por vez (sección~\ref{sec:design}).

\subsection{Regionalización de variables meteorológicas}

Sean $y_i$ las observaciones de una variable en $n$ estaciones con
coordenadas $\mathbf{u}_i$ y covariables $\mathbf{x}_i$ (elevación y distancia
a la costa). Se comparan tres familias:
\begin{itemize}
  \item \textbf{C1, estación representativa}: $\hat y(\mathbf{u}) = y_r$,
  constante en el dominio; con la variante \textbf{C1-lap}, que corrige la
  temperatura por un gradiente vertical fijo
  ($\Gamma = \SI{6,5}{\kelvin\per\kilo\metre}$).
  \item \textbf{C2, métodos convencionales}: regresión lineal global por
  mínimos cuadrados, $\hat y = \mathbf{x}^\top\hat{\boldsymbol\beta}$
  (\textbf{C2-OLS}); e interpolación por distancia inversa
  $\hat y(\mathbf{u}) = \sum_i w_i y_i/\sum_i w_i$, $w_i = d_i^{-p}$
  (\textbf{C2-IDW}).
  \item \textbf{C3, GWR}: $\hat y(\mathbf{u}) =
  \mathbf{x}^\top\hat{\boldsymbol\beta}(\mathbf{u})$ con
  $\hat{\boldsymbol\beta}(\mathbf{u}) = (\mathbf{X}^\top\mathbf{W}(\mathbf{u})
  \mathbf{X})^{-1}\mathbf{X}^\top\mathbf{W}(\mathbf{u})\mathbf{y}$, donde
  $\mathbf{W}$ contiene pesos de un núcleo bicuadrado con ancho de banda
  adaptativo ($k$ vecinos) seleccionado por AICc
  \parencite{brunsdon1996gwr,fotheringham2002gwr,gollini2015gwmodel}.
\end{itemize}
GWR reduce a la regresión global cuando el ancho de banda tiende a infinito.
Su flexibilidad tiene costos conocidos: varianza elevada de los coeficientes
locales con pocas estaciones, colinealidad local
\parencite{wheeler2005multicollinearity} y extrapolación fuera del rango de
las covariables. La pregunta de si GWR \emph{regionaliza mejor} se responde
fuera de muestra, con validación cruzada espacial que vuelva a seleccionar
el ancho de banda dentro de cada conjunto de entrenamiento
\parencite{roberts2017cv,valavi2019blockcv}; es distinta de la pregunta de
si produce \emph{mejor} $\ETa$.
